\section{视频模型的3D理解能力}

\subsection{2D生成模型的优势}

与3D生成模型形成鲜明对比的是，2D图像和视频生成模型能够很好地随数据和计算资源扩展（scale）：互联网上有数十亿张2D图像和视频可供训练，且这些模型直接生成2D像素，归纳偏置有限。随着算力不断增长，这些模型产出的视频质量越来越高——例如 V3 模型已能生成具有多样相机运动、场景运动以及高度逼真物理效果的视频。

\subsection{视频模型是否理解3D？}

一个自然的问题是：这些视频模型生成的结果是否在物理上是正确的？

Henzler 团队设计了一个实验来探究这个问题：

\begin{enumerate}
    \item 给定两张重叠区域很少的场景图像 A 和 B
    \item 使用现成的位姿估计器（如 DUSt3R）预测相机位姿——由于图像重叠少，基于视觉对应关系的估计器效果很差
    \item 使用多种视频生成模型在两张图像之间插值生成中间帧
    \item 将生成的中间帧连同原始两张图像一起输入位姿估计器
\end{enumerate}

\begin{importantbox}{关键发现：视频模型隐式理解3D}
实验结果表明，视频模型生成的中间帧显著提升了位姿估计器的精度。这意味着视频模型确实能够隐式地建模物理世界——它们生成的帧之间存在合理的3D几何关系，即使模型从未被显式训练去理解3D结构。
\end{importantbox}

\subsection{交互式视频模型}

更进一步，类似 Genie 2 的交互式视频模型已经允许用户实时导航3D世界。用户可以通过键盘输入（前进、左转、右转）来控制模型生成新帧，自由地在虚拟场景中探索不同视角。更关键的是，当用户转身回看时，场景保持了\textbf{一致性}——教室看起来与之前完全相同。

\subsection{为何仍需显式3D建模？}

既然纯2D视频模型已经如此强大，为什么还需要显式建模3D？

\begin{warningbox}{视频模型的根本瓶颈：推理成本}
虽然视频模型能随数据和计算扩展，但其推理成本也会相应上升。这使得这些模型目前无法在手机和平板等移动设备上运行。而显式3D表示（如 NeRF、3D Gaussian Splatting）一旦构建完成，渲染成本极低，完全可以在移动端实时运行。
\end{warningbox}

\subsection{两块拼图的组合}

Henzler 指出，我们其实已经拥有了两块互补的拼图：

\begin{itemize}
    \item \textbf{生成式图像/视频模型}：产出高度逼真的2D图像
    \item \textbf{NeRF/3DGS 重建方法}：将多张图像蒸馏为高效的3D表示，支持快速渲染
\end{itemize}

将二者组合——用生成模型产出多视角图像，再用3D重建方法将其蒸馏为高效3D表示——正是后续工作 CAT3D 和 B3D 的核心思路。

\subsection{本章小结}

视频模型已经展现出隐式的3D理解能力，但其推理成本限制了在移动端的部署。显式3D表示提供了低成本渲染的优势，因此「2D生成 + 3D蒸馏」成为兼顾质量与效率的最优策略。

